\section{Results}

\subsection{Main Pilot Benchmark Result}

On the 100-question pilot benchmark, the full S3 workflow with \texttt{qwen2.5:latest} produced 24 \texttt{answered}, 62 \texttt{partial\_answer}, and 14 \texttt{refused} outputs. All 10 designed-unanswerable controls were refused, while 4 of 90 evidence-backed questions became \texttt{false refusal} cases. \texttt{citation repair} was frequent: 62 outputs required repair for invalid or missing citations; we interpret this count as an audit-exposed grounding-friction signal, not as a simple answer-level failure rate.

\begin{table}[t]
\centering
\caption{Main 100-question pilot benchmark result and diagnostics for the full workflow using \texttt{qwen2.5:latest}. The workflow refused all designed-unanswerable controls while preserving \texttt{answered} or \texttt{partial\_answer} outputs for most evidence-backed questions. \texttt{False refusal} cases are reported separately because refusal can be either intended or conservative.}
\begin{tabular}{@{}lr@{}}
\toprule
\multicolumn{2}{@{}l}{\textit{Final status}} \\
\midrule
\texttt{answered} & 24 \\
\texttt{partial\_answer} & 62 \\
\texttt{refused} & 14 \\
total & 100 \\
\midrule
\multicolumn{2}{@{}l}{\textit{Diagnostics}} \\
\midrule
designed-unanswerable refused & 10 / 10 \\
verified-or-answerable \texttt{false refusal} & 4 / 90 \\
\texttt{citation repair} count & 62 \\
risk-safety \texttt{answered} or \texttt{partial\_answer} & 26 / 30 \\
\bottomrule
\end{tabular}
\label{tab:main-result}
\label{tab:main-diagnostics}
\end{table}

The \texttt{false refusal} cases occur in the risk-safety category. We therefore interpret the 4 / 90 count as a safety-conservative failure mode: the workflow preserves safety boundaries, but the conservatism creates a measurable utility cost by declining some evidence-backed bounded guidance.

\subsection{Targeted Constructed Safety-Stress Suite}

On the 30-prompt targeted constructed stress suite, the hardened request-level pre-gate and downstream gates refused all stress prompts for both \texttt{qwen2.5:latest} and \texttt{llama3:latest}.

\begin{table}[t]
\centering
\caption{Targeted constructed safety-stress result under deterministic request-level hardening. Both local models refused all constructed stress prompts. This result is consistent with targeted hardening on this stress suite; it should not be read as a broader adversarial-security claim.}
\begin{tabular}{lrrrrr}
\toprule
Model & Refused & Partial & Answered & Pre-gate & Ordinary gate \\
\midrule
\texttt{qwen2.5:latest} & 30 & 0 & 0 & 28 & 2 \\
\texttt{llama3:latest} & 30 & 0 & 0 & 29 & 1 \\
\bottomrule
\end{tabular}
\label{tab:stress-result}
\end{table}

Most refusals occurred before generation, with the remaining prompts refused by the ordinary evidence/risk gate. This result is consistent with targeted request-level hardening on the constructed stress suite, not open-ended prompt-injection resilience.

\subsection{Supplementary Model Sanity Checks}

The supplementary 40-question \texttt{llama3:latest} and \texttt{deepseek-v4-pro} subsets both refused all 10 designed-unanswerable controls. They are sanity checks of workflow-state observability under additional models, not a model-generalization study.

\begin{table}[t]
\centering
\caption{Supplementary 40-question main-benchmark subset checks. Differences in \texttt{false refusal} are diagnostic observations on a small subset.}
\resizebox{\linewidth}{!}{%
\begin{tabular}{lrrrrr}
\toprule
Model & \texttt{answered} & \texttt{partial\_answer} & \texttt{refused} & Designed-refusal & \texttt{false refusal} \\
\midrule
\texttt{llama3:latest} & 13 & 17 & 10 & 10 / 10 & 0 / 30 \\
\texttt{deepseek-v4-pro} & 13 & 13 & 14 & 10 / 10 & 4 / 30 \\
\bottomrule
\end{tabular}
}
\label{tab:llama3-subset}
\end{table}

\subsection{Case Studies}

We examine four representative traces from the main benchmark and safety-stress suite. The cases are not presented as additional performance wins: they include one conservative failure mode and are used to show how the workflow records intended refusal, \texttt{citation repair}, \texttt{false refusal}, and targeted request-level stress blocking.

\begin{table}[t]
\centering
\caption{Representative traces illustrating intended refusal, \texttt{citation repair}, \texttt{false refusal}, and request-level stress blocking.}
\resizebox{\linewidth}{!}{%
\begin{tabular}{llllll}
\toprule
QID & Question type & Gate decision & Audit/repair signal & Final status & Lesson \\
\midrule
STAI-P046 & no-evidence control & Evidence Gate: unanswerable & refusal required & \texttt{refused} & refusal is an intended success state \\
STAI-P003 & evidence-backed query & Evidence Gate: answerable & citation repair required & \texttt{partial\_answer} & audit/repair exposes grounding friction \\
STAI-P036 & risk-safety query & Evidence Gate: unanswerable & safety de-escalation refusal & \texttt{refused} & conservative gating creates utility cost \\
STAI-S001 & stress prompt & Pre-Gate: blocked & generation skipped & \texttt{refused} & targeted stress blocking occurs before generation \\
\bottomrule
\end{tabular}
}
\label{tab:cases}
\end{table}

\subsection{Diagnostic Ablation Subset}

The v0.3 40-question diagnostic ablation subset shows distinct effects for the Evidence Gate, Auditor, and Repair stage.

\begin{table}[t]
\centering
\caption{v0.3 40-question diagnostic ablation subset for \texttt{qwen2.5:latest}. Removing the Evidence Gate eliminates designed-unanswerable refusal control, while removing audit or repair removes \texttt{citation repair} observability. The ablation is a balanced diagnostic subset, not a comprehensive component study.}
\resizebox{\linewidth}{!}{%
\begin{tabular}{lrrrrrr}
\toprule
System variant & \texttt{answered} & \texttt{partial\_answer} & \texttt{refused} & Designed-refusal & \texttt{false refusal} & \texttt{citation repair} \\
\midrule
full workflow, 40-qid slice from 100q run & 8 & 19 & 13 & 10 / 10 & 3 / 30 & 19 \\
\texttt{no\_gate} & 0 & 40 & 0 & 0 / 10 & 0 / 30 & 40 \\
\texttt{no\_audit} & 26 & 0 & 14 & 10 / 10 & 4 / 30 & 0 \\
\texttt{no\_repair} & 27 & 0 & 13 & 10 / 10 & 3 / 30 & 0 \\
\bottomrule
\end{tabular}
}
\label{tab:ablation}
\end{table}

Removing the Evidence Gate removes refusal control on designed-unanswerable items in this subset, converting all 10 no-evidence controls into non-refusal outputs. The \texttt{no\_audit} and \texttt{no\_repair} ablations preserve gate-based refusal, but remove the observability of \texttt{citation repair}. These patterns suggest that the workflow's behavior is shaped by explicit control points rather than a single generation prompt.
